\section{Introduction}\label{sec:introduction}

The generalized star-height problem asks how much nested iteration is
needed to describe regular languages when complement is available.
Complement can replace some uses of iteration, but its presence does not
make ordinary star-height arguments applicable: a lower bound proved for
expressions without complement need not survive its addition. Our aim is
an absolute bound, independent of the alphabet and of the recognizing
automaton, obtained from a fixed number of operations on finite computations.

A generalized regular expression over a finite alphabet \(\Sigma\) is formed
from \(0,1\), and the letters of \(\Sigma\), using union, concatenation,
complement, and Kleene star. The constants denote the empty language and
\(\{\epsilon\}\), and a letter denotes its singleton language. Every
complement is relative to the same \(\Sigma^*\). Define the height by
\[
\begin{aligned}
 \sht(0)=\sht(1)=\sht(a)&=0,\\
 \sht(P+Q)=\sht(PQ)&=\max\{\sht(P),\sht(Q)\},\\
 \sht(\neg P)&=\sht(P),&
 \sht(P^*)&=\sht(P)+1.
\end{aligned}
\]
For a regular language \(L\subseteq\Sigma^*\), let \(h_\Sigma(L)\) be the
minimum of these heights among expressions denoting \(L\). Here regular
means accepted by a deterministic finite automaton (DFA), with a finite
state set, letter transitions, an initial state, and accepting states.
There is no restriction on expression length. The uniform-boundedness question asks
whether one integer bounds \(h_\Sigma(L)\) for every finite \(\Sigma\) and
every regular \(L\) over it. These conventions agree with the standard
Boolean expression model in \cite[\S1]{PST1992}.

\begin{theorem}\label{thm:main}
For every finite alphabet \(\Sigma\) and every regular language
\(L\subseteq\Sigma^*\),
\[
                         h_\Sigma(L)\le 13.
\]
\end{theorem}

This gives an affirmative answer to uniform boundedness. It does not settle
whether height one always suffices. Straubing explicitly distinguishes
these two questions~\cite[p.~537]{Straubing2002}: a language requiring
height greater than one would rule out the bound one, but would not rule
out a larger absolute bound. Computing the minimum generalized height of
a given language is a further question. The construction below supplies
a uniform upper bound; it does not determine that minimum.

The complete construction below uses prediction trials and full decision
histories. Two separate companion articles prove the smaller bounds
four~\cite[Theorem~1.1]{OpenAIHeight4} and
three~\cite[Theorem~1.1]{OpenAIHeight3}, respectively by multiscale
periodic transport and by arbitrary-function shift tables with a
whole-stencil decoder. Each article proves its own needed local statements.
Here the value thirteen records the depth of the prediction/history method;
the stronger numerical bounds do not replace any step of its proof.

\subsection*{Historical and algebraic context}

Eggan related ordinary star height to transition graphs~\cite{Eggan1963},
and Dejean and Sch\"utzenberger gave ordinary-height examples over a
binary alphabet~\cite{DejeanSchutzenberger1966}. Their expressions omit
complement, so their lower bounds do not transfer to the generalized
model. The effect of complement already appears at height zero:
\(\Sigma^*=\neg0\) is star-free. Sch\"utzenberger's 1965 theorem
characterizes the star-free languages as those recognized by finite
aperiodic monoids~\cite{Schutzenberger1965}. Here a finite monoid is
aperiodic if each element \(x\) satisfies \(x^n=x^{n+1}\) for some \(n\).

This characterization makes finite monoids a natural setting for the
problem. A word acts on the states of a DFA, and the induced
transformations form a finite monoid under composition. Height-one
results cover several algebraic classes. Henneman's
result for finite abelian groups is recalled and proved by Pin,
Straubing, and Th\'erien~\cite[Theorem~3.4]{PST1992}; their own
Theorem~7.3 establishes height at most one for languages recognized by
finite nilpotent groups of class two. Bourne and Ru\v{s}kuc subsequently
proved the same bound for languages recognized by finite Rees
zero-matrix semigroups over abelian
groups~\cite[Theorem~3.6]{BourneRuskuc2016}. These results treat specified
recognizing classes; the construction here applies to every finite
transition monoid.

The alphabet and morphism conventions are essential to this comparison.
Pin, Straubing, and Th\'erien prove that star-free injective substitutions
and inverse alphabetic morphisms do not increase generalized
height~\cite[Theorem~4.6 and Corollary~4.7]{PST1992}. Their Theorem~8.1
also represents every regular language as an inverse image, under a
general morphism, of a language of ordinary height at most one. The
inverse-image closure above is restricted to alphabetic morphisms, so it
does not turn this representation into a height-one theorem. Our proof
constructs every expression directly over the original alphabet.

Place and Zeitoun's 2017 account distinguishes the ordinary and
generalized problems and records that no language of generalized height
greater than one was then known~\cite[\S1]{PlaceZeitoun2017}.
Cotumaccio's February 2025 discussion still identifies generalized star
height as an open problem connected with automata and
compression~\cite[p.~46]{Cotumaccio2025}. The theorem above addresses
uniform boundedness; the two stronger companion bounds stated earlier
also do not settle whether height one suffices.

The local-marker ingredient comes from a different tradition. Krieger's
separated-marker construction~\cite[Lemma~2]{Krieger1982} places markers
apart while controlling markerless regions by periodicity. Meyerovitch
presents a clopen-marker form and its finite merging
argument~\cite[\S3, Lemmas~3.2--3.4]{Meyerovitch2025}. Section~\ref{sec:geometry}
proves the precise finite-window construction and periodicity estimates
needed here. The historical results provide context; every local lemma
used to obtain the bound thirteen is proved below.

\subsection*{The mechanism of the proof}

Fix the finite monoid of a recognizing DFA. We represent its action on a
finite vector space over \(\mathbf F_2\). A word will be parsed into
nonempty pieces, called episodes, and one final remainder. The episodes
form a prefix code: no episode is a proper prefix of another. Consequently
their parsing is unique. Each episode carries a total finite-state update.
A graph language specifies the initial and final states of the composed
update on an episode sequence.

The basic step splits an episode into a prefix and a suffix. The prefix
checks the incoming state and the suffix checks the outgoing state. Each
side may choose its own finite metadata. For soundness, every pair of
successful choices must describe the actual episode update and consume
exactly one episode. For availability, every designated episode and every
input state must admit some successful cut. If graph languages for the
remaining episodes are already known, one additional star combines their
sequences with the successful splits. Lemma~\ref{lem:split} states this
identity and its exact height cost.

Three constructions make a bounded number of these steps sufficient.
\begin{enumerate}[leftmargin=*,label=\textup{\arabic*.}]
\item \emph{An affine computation tolerates a controlled disagreement.}
A map \(v\mapsto vA+b\) has linear part \(A\), which is recovered
from its values at \(v\) and at zero. Our finite clock ensures that,
at a given total length residue, independently successful tags either
agree or have one fixed off-diagonal pair. In the latter case the
claimed input and output fit one affine correction with the required
linear part. The clock also bounds how many ordered pairs of candidate
cuts are spoiled, independently of the size of their finite roster.
\item \emph{Prediction and periodicity handle two large episode classes.}
A trial predicts the monoid products on a fixed subdivision and offers
one block of cuts for each input--output vector pair. Partial prefix
and suffix checks always fit a deterministic affine trial update;
a perfect prediction supplies a cut for every input. Local markers
select the trial boundaries. When a boundary window is markerless,
periodicity supplies a second family of cuts with a constant monoid
prefix product and constant context. These are the two outer splits.
\item \emph{Five residual witnesses permit exact time decoding.}
Failure of the two availability conditions has one of five explicit
types. At each type we place many copies of every required cut kind.
A suffix checks all candidate origins and identifies the unique one
having the current witness after later witnesses have been excluded.
A bound on ambiguous ordered pairs leaves a usable copy of each kind.
One aligned split then uses a perfect trial. If the actual full decision
history has no perfect trial, a second split recovers a measurement
from that history. Both splits compute the same affine update on a
larger space; its linear part encodes the target update and resets an
auxiliary monoid register, so the encoding composes across episodes.
The argument counts all leaves, including those that no actual prefix
can realize; it never cancels a monoid factor.
\end{enumerate}

The finite objects in this description can be extremely large. The
construction first fixes their state sets and event kinds, then their
copy counts, clocks, marker scales, and positions in a compatible order.
Only the number of nested split steps matters for height: there are two
outer steps and two steps for each of five residual types. Their component
languages have height at most one because the episode-end rule restricts
unbounded words to a fixed periodic middle between bounded ends. The
final compilation includes the incomplete remainder and yields the
thirteen-height count without an alphabet substitution.

\paragraph{Organization.}
Section~\ref{sec:algebra} develops finite computation, the split identity,
and prediction trials. Section~\ref{sec:clock} proves the finite clock.
Section~\ref{sec:geometry} constructs local markers and exact episode-end
guards. Section~\ref{sec:outer} gives the two outer splits and their five
residual witnesses; Section~\ref{sec:aligned} eliminates those witnesses
with full-history transmission. Section~\ref{sec:parameters} makes all
finite choices compatible. Section~\ref{sec:compilation} gives explicit
original-alphabet expressions and completes the proof.
